\capitulo{2}{Objetivos del proyecto}

El propósito principal de este proyecto es el desarrollo, y ejecución de un sistema de software modular destinado a registrar, analizar y mostrar en automático las interacciones del usuario en entornos de modelado tridimensional, utilizando Blender como base para tal desarrollo.

La función central de este programa es proporcionar un sistema fuerte que capture y organize información del proceso de diseño de forma automática y transparente, de forma que permita un análisis estadístico posterior. Por tanto, el objetivo es conseguir métricas objetivas sobre el comportamiento del usuario, el nivel de uso de herramientas y la transformación geométrica de las escenas, aportando un mayor conocimiento del proceso de modelado tridimensional~\cite{jankowski2014,lavioila2017}.

La solución desarrollada se enfoca en capturar información útil sobre el proceso de modelado sin afectar el rendimiento del entorno o el flujo de trabajo del usuario. Con esto, es posible analizar los conjuntos de datos generados para investigar cómo se construyen los modelos tridimensionales, cuáles operaciones son más comunes y qué acontecimientos ocurren dentro del entorno virtual~\cite{bowman2004}.

El proyecto considera todo el proceso de recopilación y conservación de datos, así como su posterior procesamiento analítico y visualización. El sistema dispone de técnicas capaces de explorar la información recolectada desde múltiples perspectivas: el seguimiento del tiempo, el análisis de la evolución de la malla geométrica, el descubrimiento de eventos de modelado por requerimiento y la representación gráfica de métricas directamente en el mismo espacio tridimensional. Este modelo coincide con las estrategias actuales vinculadas al análisis de la interacción en entornos virtuales y a los sistemas avanzados de modelado asistido~\cite{malik2023}.

\section{Objetivos específicos}

Para alcanzar el objetivo general planteado, se definen los siguientes objetivos específicos:

\begin{itemize}

\item \textbf{Diseño de un núcleo de adquisición no intrusivo}: Desarrollar un sistema de captura de datos en segundo plano que funcione de manera transparente y sin latencia perceptible durante las sesiones de modelado 3D.

\item \textbf{Implementación de una arquitectura modular de extensión}: Programar un conjunto de \textit{add-ons} para Blender, integrados en su API oficial, capaces de extraer información temporal, espacial y estructural de la escena activa y de los objetos que la componen.

\item \textbf{Automatización del análisis topológico complejo}: Desarrollar algoritmos específicos para la detección automática de eventos analíticos avanzados, tales como duplicaciones de vértices, fusiones geométricas, conmutaciones de modo, cuantificación de islas de coordenadas UV y cálculo de distancias de similitud geométrica.

\item \textbf{Garantía de persistencia externa y estructurada}: Diseñar un mecanismo de almacenamiento eficiente basado en archivos de texto plano estructurados (formatos de persistencia CSV), asegurando la integridad, portabilidad y trazabilidad de los datos entre diferentes sesiones de trabajo de forma independiente a la plataforma.

\item \textbf{Establecimiento de un marco ético de captura}: Incorporar un módulo de consentimiento y configuración de privacidad en la interfaz de usuario que garantice una recopilación de datos ética, informada y transparente.

\item \textbf{Desarrollo de un motor de visualización interactivo}: Diseñar e implementar componentes de representación gráfica y procedural (gráficos estadísticos, mapas de dispersión y diagramas de radar translúcidos) integrados de forma nativa en el visor 3D de Blender para facilitar la interpretación del comportamiento del usuario.

\item \textbf{Definición de un sistema de métricas interpretables}: Organizar los indicadores calculados por el sistema mediante una nomenclatura estable, incluyendo métricas estadísticas A1--A17, métricas geométricas, métricas UV y métricas de actividad operacional, de modo que puedan relacionarse de forma clara con los datos capturados.

\item \textbf{Optimización de la interfaz analítica}: Refinar la presentación de resultados mediante bloques compactos, agrupaciones de métricas relacionadas, textos visibles breves y ayudas contextuales, reduciendo la sobrecarga visual sin eliminar información técnica relevante.

\item \textbf{Validación funcional preliminar del sistema}: Comprobar la instalación, activación, captura, exportación, análisis y visualización mediante pruebas automatizadas sobre la lógica desacoplada y pruebas manuales controladas dentro de Blender. La evaluación experimental con usuarios finales queda fuera del alcance cerrado de esta versión y se plantea como trabajo futuro.

\item \textbf{Preparación para análisis conductuales futuros}: Definir métricas y estructuras de datos que puedan apoyar estudios posteriores sobre estrategias de modelado. La identificación de diferencias significativas entre usuarios expertos y novatos requeriría una muestra experimental, un protocolo de evaluación y un análisis estadístico específicos.

\end{itemize}

Estos objetivos específicos permiten estructurar el desarrollo del proyecto en fases metodológicas bien definidas, abarcando desde la ingeniería de requisitos y la implementación del software hasta el análisis preliminar de la información y la preparación de futuras validaciones experimentales en entornos reales de aprendizaje.
